% GENERATED FILE. Edit canonical lessons, then run npm run generate:books.
% canonical-source-hash: fnv1a64:40aa7e30fd05ce86
% canonical-lessons: RU-C283-utverzhdat, RU-C283-upominat, RU-C283-utochnyat, RU-C283-vyyasnyat, RU-C283-preduprezhdat

\chapter{Claiming and checking}
\label{ch:ru-claiming-and-checking}
\hlchaptermodality{283}

\begin{chapteropening}
\textbf{By the end of this chapter:} \emph{I can say that someone claimed, mentioned, checked, found out or warned about something.}

Complete the last lesson of chapter 283: I can say that someone claimed, mentioned, checked, found out or warned about something.
\end{chapteropening}

\section[\texorpdfstring{\ru{утверждать} (utverzhdát)}{утверждать (utverzhdát)}]{\ru{утверждать} (utverzhdát) — to claim, to assert}

\label{lesson:RU-C283-utverzhdat}



\begin{quote}
Before the new one: say the Russian for a detail, then the Russian for a fact.
\end{quote}

\subsection*{You'll want to know: \ru{утверждать}}
\textbf{\ru{утверждать}} (\emph{utverzhdát}) — \textquotedblleft{}to claim, to assert\textquotedblright{}.

\begin{quote}
\ru{Он} \ru{утверждает}, \ru{что} \ru{видел} \ru{медведя}.
\end{quote}

\begin{quote}
\emph{On utverzhdáyet, shto vídel medvédya.}
\end{quote}

\begin{quote}
He claims that he saw a bear.
\end{quote}

In this sense imperfective. The perfective \emph{\ru{утвердить}} (\emph{utverdít}) means \textquotedblleft{}to approve\textquotedblright{}: \emph{\ru{утвердить} \ru{план}}.

\begin{cousinweb}
\emph{\ru{утверждать}} is built on \emph{\ru{у}-} and the root of \emph{\ru{твёрдый}} \textquotedblleft{}firm\textquotedblright{}, Proto-Slavic \emph{tvĭrdŭ}, in its Church Slavonic shape \emph{-\ru{твержд}-}; EtymDB lists no parent for the verb.
\end{cousinweb}

\begin{grammarlens}[title={What you've built}]
One more word for telling, explaining and describing at B1.
\end{grammarlens}

\subsection*{Your turn}
\begin{itemize}
\raggedright
  \item \emph{Say it:} \emph{utverzhdát}
  \item \emph{Say it:} \emph{utverzhdát}, once more
  \item \emph{Say it:} a sentence of your own with \emph{utverzhdát}
  \item \emph{Recall:} say the Russian for a detail, then the Russian for a fact, then say \emph{utverzhdát} again
\end{itemize}

\begin{tcolorbox}[breakable,colback=teal!4,colframe=teal!35!black,title={Before you move on}]
What does \ru{утверждать} mean? (\textquotedblleft{}To claim, to assert\textquotedblright{}.) Say it once more.
\end{tcolorbox}
\section[\texorpdfstring{\ru{упоминать} (upominát)}{упоминать (upominát)}]{\ru{упоминать} (upominát) — to mention}

\label{lesson:RU-C283-upominat}



\begin{quote}
Before the new one: say the Russian for a fact, then the Russian for to claim.
\end{quote}

\subsection*{You'll want to know: \ru{упоминать}}
\textbf{\ru{упоминать}} (\emph{upominát}) — \textquotedblleft{}to mention\textquotedblright{}.

\begin{quote}
\ru{В} \ru{письме} \ru{она} \ru{упомянула} \ru{твоё} \ru{имя}.
\end{quote}

\begin{quote}
\emph{V pismé oná upomyanúla tvoyó ímya.}
\end{quote}

\begin{quote}
In the letter she mentioned your name.
\end{quote}

Imperfective; perfective \emph{\ru{упомянуть}} (\emph{upomyanút}). It takes the accusative or \emph{\ru{о}} + prepositional.

\begin{cousinweb}
\emph{\ru{упоминать}} is built on \emph{\ru{у}-} and the Slavic root of \emph{\ru{помнить}} \textquotedblleft{}to remember\textquotedblright{} and \emph{\ru{память}} \textquotedblleft{}memory\textquotedblright{}: to bring to mind. EtymDB lists no parent.
\end{cousinweb}

\begin{grammarlens}[title={What you've built}]
One more word for telling, explaining and describing at B1.
\end{grammarlens}

\subsection*{Your turn}
\begin{itemize}
\raggedright
  \item \emph{Say it:} \emph{upominát}
  \item \emph{Say it:} \emph{upominát}, once more
  \item \emph{Say it:} a sentence of your own with \emph{upominát}
  \item \emph{Recall:} say the Russian for a fact, then the Russian for to claim, then say \emph{upominát} again
\end{itemize}

\begin{tcolorbox}[breakable,colback=teal!4,colframe=teal!35!black,title={Before you move on}]
What does \ru{упоминать} mean? (\textquotedblleft{}To mention\textquotedblright{}.) Say it once more.
\end{tcolorbox}
\section[\texorpdfstring{\ru{уточнять} (utochnyát)}{уточнять (utochnyát)}]{\ru{уточнять} (utochnyát) — to clarify, to check (details)}

\label{lesson:RU-C283-utochnyat}



\begin{quote}
Before the new one: say the Russian for to claim, then the Russian for to mention.
\end{quote}

\subsection*{You'll want to know: \ru{уточнять}}
\textbf{\ru{уточнять}} (\emph{utochnyát}) — \textquotedblleft{}to clarify, to check (details)\textquotedblright{}.

\begin{quote}
\ru{Я} \ru{позвонил}, \ru{чтобы} \ru{уточнить} \ru{время}.
\end{quote}

\begin{quote}
\emph{Ya pozvoníl, shtóby utochnít vrémya.}
\end{quote}

\begin{quote}
I phoned to check the time.
\end{quote}

Imperfective; perfective \emph{\ru{уточнить}} (\emph{utochnít}).

\begin{cousinweb}
\emph{\ru{уточнять}} is built on \emph{\ru{у}-} and \emph{\ru{точный}} \textquotedblleft{}exact\textquotedblright{}, which EtymDB derives from Church Slavonic \emph{tŭčĭnŭ}; it is usually traced further to the root of \emph{\ru{ткнуть}} \textquotedblleft{}to poke\textquotedblright{}, so \textquotedblleft{}hitting the spot\textquotedblright{}.
\end{cousinweb}

\begin{grammarlens}[title={What you've built}]
One more word for telling, explaining and describing at B1.
\end{grammarlens}

\subsection*{Your turn}
\begin{itemize}
\raggedright
  \item \emph{Say it:} \emph{utochnyát}
  \item \emph{Say it:} \emph{utochnyát}, once more
  \item \emph{Say it:} a sentence of your own with \emph{utochnyát}
  \item \emph{Recall:} say the Russian for to claim, then the Russian for to mention, then say \emph{utochnyát} again
\end{itemize}

\begin{tcolorbox}[breakable,colback=teal!4,colframe=teal!35!black,title={Before you move on}]
What does \ru{уточнять} mean? (\textquotedblleft{}To clarify, to check (details)\textquotedblright{}.) Say it once more.
\end{tcolorbox}
\section[\texorpdfstring{\ru{выяснять} (vyyasnyát)}{выяснять (vyyasnyát)}]{\ru{выяснять} (vyyasnyát) — to find out, to clear up}

\label{lesson:RU-C283-vyyasnyat}



\begin{quote}
Before the new one: say the Russian for to mention, then the Russian for to clarify.
\end{quote}

\subsection*{You'll want to know: \ru{выяснять}}
\textbf{\ru{выяснять}} (\emph{vyyasnyát}) — \textquotedblleft{}to find out, to clear up\textquotedblright{}.

\begin{quote}
\ru{Мы} \ru{выяснили}, \ru{кто} \ru{взял} \ru{ключ}.
\end{quote}

\begin{quote}
\emph{My výyasnili, kto vzyal klyuch.}
\end{quote}

\begin{quote}
We found out who had taken the key.
\end{quote}

Imperfective; perfective \emph{\ru{выяснить}} (\emph{výyasnit}), stressed on \emph{\ru{вы}-}.

\begin{cousinweb}
\emph{\ru{выяснять}} is built on \emph{\ru{вы}-} \textquotedblleft{}out\textquotedblright{} and \emph{\ru{ясный}} \textquotedblleft{}clear, bright\textquotedblright{}, inherited from Proto-Slavic \emph{ěsnŭ} \textquotedblleft{}clear\textquotedblright{} (EtymDB): to bring out into the light.
\end{cousinweb}

\begin{grammarlens}[title={What you've built}]
One more word for telling, explaining and describing at B1.
\end{grammarlens}

\subsection*{Your turn}
\begin{itemize}
\raggedright
  \item \emph{Say it:} \emph{vyyasnyát}
  \item \emph{Say it:} \emph{vyyasnyát}, once more
  \item \emph{Say it:} a sentence of your own with \emph{vyyasnyát}
  \item \emph{Recall:} say the Russian for to mention, then the Russian for to clarify, then say \emph{vyyasnyát} again
\end{itemize}

\begin{tcolorbox}[breakable,colback=teal!4,colframe=teal!35!black,title={Before you move on}]
What does \ru{выяснять} mean? (\textquotedblleft{}To find out, to clear up\textquotedblright{}.) Say it once more.
\end{tcolorbox}
\section[\texorpdfstring{\ru{предупреждать} (preduprezhdát)}{предупреждать (preduprezhdát)}]{\ru{предупреждать} (preduprezhdát) — to warn}

\label{lesson:RU-C283-preduprezhdat}



\begin{quote}
Before the new one: say the Russian for to clarify, then the Russian for to find out.
\end{quote}

\subsection*{You'll want to know: \ru{предупреждать}}
\textbf{\ru{предупреждать}} (\emph{preduprezhdát}) — \textquotedblleft{}to warn\textquotedblright{}.

\begin{quote}
\ru{Мать} \ru{предупредила} \ru{детей} \ru{о} \ru{грозе}.
\end{quote}

\begin{quote}
\emph{Mat predupredíla detéy o grozé.}
\end{quote}

\begin{quote}
Mother warned the children about the storm.
\end{quote}

Imperfective; perfective \emph{\ru{предупредить}} (\emph{predupredít}), \emph{\ru{я} \ru{предупрежу}}. The person is in the accusative, the danger with \emph{\ru{о}}.

\begin{cousinweb}
\emph{\ru{предупреждать}} is built on \emph{\ru{пред}-} \textquotedblleft{}before\textquotedblright{} and \emph{\ru{упредить}} \textquotedblleft{}to get ahead of\textquotedblright{}, in the Church Slavonic shape of the root seen in native \emph{\ru{перед}}: to tell before it happens. EtymDB lists no parent.
\end{cousinweb}

\begin{grammarlens}[title={What you've built}]
One more word for telling, explaining and describing at B1.
\end{grammarlens}

\subsection*{Your turn}
\begin{itemize}
\raggedright
  \item \emph{Say it:} \emph{preduprezhdát}
  \item \emph{Say it:} \emph{preduprezhdát}, once more
  \item \emph{Say it:} a sentence of your own with \emph{preduprezhdát}
  \item \emph{Recall:} say the Russian for to clarify, then the Russian for to find out, then say \emph{preduprezhdát} again
\end{itemize}

\begin{tcolorbox}[breakable,colback=teal!4,colframe=teal!35!black,title={Before you move on}]
What does \ru{предупреждать} mean? (\textquotedblleft{}To warn\textquotedblright{}.) Say it once more.
\end{tcolorbox}
